\documentclass[11pt,a4paper]{article}

\usepackage[T1]{fontenc}
\usepackage[utf8]{inputenc}
\usepackage{lmodern}
\usepackage{microtype}
\usepackage{geometry}
\geometry{margin=1in}

\usepackage{amsmath}
\usepackage{amssymb}
\usepackage{hyperref}
\hypersetup{
  colorlinks=true,
  linkcolor=blue!50!black,
  urlcolor=blue!50!black,
  citecolor=blue!50!black
}

\title{Adaptive Sparse Grid Toolkit\\Theory Notes}
\author{}
\date{\today}

\begin{document}
\maketitle

\begin{abstract}
This document summarizes the theoretical background used by the toolkit:
Lagrange interpolation on tensor grids, the Smolyak combination technique,
and the computation of weighted quadrature weights for statistical analysis.
\end{abstract}

\tableofcontents
\newpage

\section{Tensor‑grid interpolation}

Let $y = (y_1,\ldots,y_N)$ be the input vector and let
$\{y^{(j)}_{m(i)}\}$ denote the nodes of a tensor grid associated with a multi‑index $i$.
The tensor‑grid interpolant of a scalar QoI $f(y)$ is
\[
U_i(y) := \sum_{j \le m(i)} f\!\left(y^{(j)}_{m(i)}\right)\,L^{(j)}_{m(i)}(y),
\qquad
L^{(j)}_{m(i)}(y) = \prod_{n=1}^{N} \ell^{(j_n)}_{n,m(i_n)}(y_n),
\]
where $\ell^{(j_n)}_{n,m(i_n)}$ are the 1D Lagrange basis polynomials on the
knot set of dimension $n$.

\section{Sparse‑grid interpolation (Smolyak combination)}

The sparse‑grid interpolant is formed by the combination technique:
\[
U_{\mathcal{I}}(y) = \sum_{i \in \mathcal{I}} c_i\,U_i(y),
\]
where $\mathcal{I}$ is the multi‑index set (total‑degree in this toolkit)
and $c_i$ are the Smolyak coefficients.

\section{Weighted quadrature from Lagrange bases}

Let $\rho(y)$ be the joint PDF, with per‑dimension factors $\rho_n(y_n)$ on
$\Gamma = \Gamma_1 \times \cdots \times \Gamma_N$. The weighted integral of $f$ is
approximated by integrating the tensor‑grid interpolant:
\begin{align*}
\int_{\Gamma} f(y)\rho(y)\,dy
&\approx \int_{\Gamma} U_i(y)\rho(y)\,dy \\
&= \sum_{j \le m(i)} f\!\left(y^{(j)}_{m(i)}\right)
\prod_{n=1}^{N} \left(\int_{\Gamma_n} \ell^{(j_n)}_{n,m(i_n)}(y_n)\rho_n(y_n)\,dy_n\right).
\end{align*}

Define the 1D Lagrange quadrature weights
\[
\omega^{(j_n)}_{n,m(i_n)} := \int_{\Gamma_n} \ell^{(j_n)}_{n,m(i_n)}(y_n)\rho_n(y_n)\,dy_n,
\]
then the tensor‑product weights are
\[\omega^{(j)}_{m(i)} := \prod_{n=1}^{N} \omega^{(j_n)}_{n,m(i_n)}.\]
The tensor‑grid quadrature is
\[
Q_i \approx \sum_{j \le m(i)} f\!\left(y^{(j)}_{m(i)}\right)\,\omega^{(j)}_{m(i)}.
\]

\section{Sparse‑grid quadrature (combination technique)}

Applying the Smolyak combination to quadrature yields
\[
\int_{\Gamma} f(y)\rho(y)\,dy
\approx
\sum_{i \in \mathcal{I}} c_i\,Q_i.
\]
In implementation, weights are computed on the union of sparse‑grid nodes,
accumulating tensor‑grid contributions and combining them with $c_i$.
Weights are normalized by default so that $\sum_k w_k = 1$, enabling direct
use as expectation weights.

\section{Numerical evaluation of 1D weights}

Each 1D weight is evaluated by Gauss–Legendre quadrature on $[a,b]$:
\[
\omega_j = \int_a^b \ell_j(x)\rho(x)\,dx
\approx \sum_{q=1}^Q w_q\,\rho(x_q)\,\ell_j(x_q),
\qquad Q = \texttt{quad\_order}.
\]
The quadrature points $x_q$ and weights $w_q$ are mapped from $[-1,1]$ to $[a,b]$.

\section{Lagrange‑basis coefficient export}

For each tensor grid $i$ and node $y^{(j)}_{m(i)}$, the toolkit exports the
combined coefficient
\[
\texttt{combined\_coeff} = c_i \cdot f\!\left(y^{(j)}_{m(i)}\right),
\]
which provides an explicit representation of the sparse‑grid interpolant as a
weighted sum of tensor‑product Lagrange basis functions.

\end{document}

